\documentclass[12pt,a4paper]{article}

\usepackage[utf8]{inputenc}
\usepackage[T1]{fontenc}
\usepackage[english]{babel}
\usepackage{lmodern}
\usepackage{microtype}
\usepackage[a4paper,margin=2.5cm]{geometry}
\usepackage{setspace}
\onehalfspacing

\usepackage{amsmath,amssymb,amsfonts,mathtools}
\usepackage{amsthm}
\usepackage{booktabs}
\usepackage{xcolor}
\usepackage[hidelinks]{hyperref}
\usepackage[round,authoryear]{natbib}

\newtheorem{proposition}{Proposition}
\newtheorem{remark}{Remark}

\newcommand{\E}{\mathbb{E}}
\newcommand{\tY}{\widetilde{Y}}
\newcommand{\tD}{\widetilde{D}}
\newcommand{\tZ}{\widetilde{Z}}

\title{\textbf{PLR-IV versus LATE Estimators:\\A Comparison of Estimands}}
\author{}
\date{}

\begin{document}
\maketitle
\thispagestyle{empty}

\section*{Overview}

This note clarifies what the Partially Linear IV (PLR-IV) estimator targets and how it
differs from LATE estimators such as Wald-AIPW and kappa-weighting.  The distinction
matters for comparing estimates across the two families: they solve different moment
conditions and are consistent for different population parameters.

% ------------------------------------------------------------------
\section{The LATE Target}
% ------------------------------------------------------------------

Under the standard IV assumptions---conditional instrument independence, exclusion
restriction, instrument relevance, overlap, and monotonicity---the local average treatment
effect is
\begin{align}
  \tau_{\mathrm{LATE}}
  \;=\;
  \E\!\left[Y_i(1)-Y_i(0)\;\middle|\;D_i(1)>D_i(0)\right].
  \label{eq:late}
\end{align}
This is the average treatment effect for \emph{compliers}, i.e.\ those whose treatment
status is switched by the instrument.  Both the Wald-AIPW estimator and kappa-based
weighting estimators \citep{abadie2003semiparametric} target $\tau_{\mathrm{LATE}}$
directly, relying on Abadie's (2003) result that any feature of the joint complier
distribution is identified through the kappa weight
\begin{align}
  \kappa_i
  \;=\;
  1
  -\frac{D_i(1-Z_i)}{1-p(X_i)}
  -\frac{(1-D_i)Z_i}{p(X_i)},
  \qquad
  p(X_i)=\E[Z_i\mid X_i],
\end{align}
so that $\tau_{\mathrm{LATE}} = \E[\kappa_i Y_i(1)]/\E[\kappa_i] -
\E[\kappa_i Y_i(0)]/\E[\kappa_i]$.  The key point is that the LATE is a
\emph{subpopulation average}: it averages heterogeneous individual effects over the
complier group.

% ------------------------------------------------------------------
\section{The PLR-IV Estimand}
% ------------------------------------------------------------------

The PLR-IV estimator is built on a structural model rather than a subpopulation average.
The outcome is assumed to follow a \emph{partially linear} equation with a
\emph{constant} treatment coefficient:
\begin{align}
  Y_i = D_i\,\theta_0 + g_0(X_i) + U_i,
  \qquad
  \E[U_i\mid X_i,Z_i]=0.
  \label{eq:pliv}
\end{align}
The function $g_0(X_i)$ is left unrestricted (estimated nonparametrically), so the model
is partially linear.  The instrument enters through the conditional exogeneity condition
$\E[U_i\mid X_i,Z_i]=0$, which replaces the endogeneity assumption on $D_i$.

Define the nuisance functions
\[
  \ell_0(X_i)=\E[Y_i\mid X_i],\quad
  r_0(X_i)=\E[D_i\mid X_i],\quad
  m_0(X_i)=\E[Z_i\mid X_i],
\]
and the corresponding residuals $\tY_i=Y_i-\ell_0(X_i)$,
$\tD_i=D_i-r_0(X_i)$, $\tZ_i=Z_i-m_0(X_i)$.
Substituting \eqref{eq:pliv} into the moment condition $\E[U_i\,\tZ_i]=0$ yields the
Neyman-orthogonal score of \citet{chernozhukov2018double}:
\begin{align}
  \psi(W_i;\theta,\eta)
  \;=\;
  \bigl(\tY_i - \theta\,\tD_i\bigr)\tZ_i
  \;\stackrel{!}{=}\; 0.
  \label{eq:pliv_score}
\end{align}
Solving for $\theta_0$ gives the population ratio
\begin{align}
  \theta_0
  \;=\;
  \frac{\E[\tZ_i\,\tY_i]}{\E[\tZ_i\,\tD_i]},
  \label{eq:pliv_ratio}
\end{align}
provided $\E[\tZ_i\,\tD_i]\neq 0$ (residualized relevance condition).  The DML
estimator replaces $\ell_0,r_0,m_0$ with cross-fitted ML estimates $\hat\ell,\hat
r,\hat m$:
\begin{align}
  \hat\theta_{\mathrm{PLR\text{-}IV}}
  \;=\;
  \frac{
    \sum_{i=1}^{N}(Z_i-\hat m(X_i))(Y_i-\hat\ell(X_i))
  }{
    \sum_{i=1}^{N}(Z_i-\hat m(X_i))(D_i-\hat r(X_i))
  }.
  \label{eq:pliv_dml}
\end{align}
Under product-rate conditions
$\|\hat m - m_0\|_{P,2}(\|\hat r-r_0\|_{P,2}+\|\hat\ell-\ell_0\|_{P,2})=o(N^{-1/2})$,
cross-fitting, and Neyman orthogonality, $\hat\theta_{\mathrm{PLR\text{-}IV}}$ is
root-$N$ consistent and asymptotically normal for $\theta_0$.

% ------------------------------------------------------------------
\section{Why They Differ}
% ------------------------------------------------------------------

The ratio in \eqref{eq:pliv_ratio} has a Wald-type form, but it is \emph{not} the LATE
Wald estimand.  The following proposition makes the distinction precise.

\begin{proposition}[PLR-IV versus LATE]
\label{prop:pliv_late}
Suppose both the PLIV model \eqref{eq:pliv} and the standard IV assumptions hold.
Then
\begin{enumerate}
  \item[\normalfont(i)] $\theta_0$ as defined in \eqref{eq:pliv_ratio} is the unique
        scalar satisfying the moment condition \eqref{eq:pliv_score}.  It is the
        structural coefficient in the constant-effect linear model.
  \item[\normalfont(ii)] $\tau_{\mathrm{LATE}}$ as defined in \eqref{eq:late} is the
        average treatment effect for compliers.  It does not require effect homogeneity.
  \item[\normalfont(iii)] $\theta_0 = \tau_{\mathrm{LATE}}$ if and only if treatment
        effects are homogeneous, i.e.\ $Y_i(1)-Y_i(0)=c$ a.s.\ for some constant
        $c\in\mathbb{R}$.
\end{enumerate}
\end{proposition}

\begin{proof}
\textit{Part (i).}  By construction, \eqref{eq:pliv_ratio} is the unique solution to
$\E[(\tY_i-\theta\tD_i)\tZ_i]=0$ when $\E[\tZ_i\tD_i]\neq 0$.

\textit{Part (ii).}  Under the standard IV assumptions, Imbens and Angrist (1994) show
that the standard Wald ratio $\E[Y_i\mid Z_i=1]-\E[Y_i\mid Z_i=0]$ divided by the
first-stage $\E[D_i\mid Z_i=1]-\E[D_i\mid Z_i=0]$ equals $\tau_{\mathrm{LATE}}$.  This
identification result holds without any constant-effect assumption.

\textit{Part (iii).}  If $Y_i(1)-Y_i(0)=c$ for all $i$, then the structural model
\eqref{eq:pliv} holds with $\theta_0=c=\tau_{\mathrm{ATE}}=\tau_{\mathrm{LATE}}$.
Conversely, if treatment effects are heterogeneous, $\theta_0$ is the unique solution to
\eqref{eq:pliv_score}, which need not equal the complier average $\tau_{\mathrm{LATE}}$.
To see this, expand the PLIV moment condition at the true parameter.  Under
heterogeneous effects write $Y_i=D_i\tau_i+g_0(X_i)+\varepsilon_i$.  Then
$\tY_i-\theta_0\tD_i=\tD_i(\tau_i-\theta_0)+\tilde\varepsilon_i$ and the moment
condition $\E[(\tau_i-\theta_0)\tD_i\tZ_i]=0$ pins down $\theta_0$ as a
covariance-weighted average of individual effects $\tau_i$, with weights proportional to
$\tD_i\tZ_i$, not the kappa-complier weights.  This weighted average does not, in
general, equal the unweighted complier mean $\tau_{\mathrm{LATE}}$.
\end{proof}

% ------------------------------------------------------------------
\section{Implications for Comparison}
% ------------------------------------------------------------------

Proposition~\ref{prop:pliv_late} has two direct consequences for the empirical analysis.

\medskip\noindent\textbf{Different targets under heterogeneity.}
In settings with heterogeneous treatment effects---which both the Card (1995) and
Kitagawa (2015) specifications implicitly allow---PLR-IV and LATE estimators are
consistent for distinct parameters.  Numerical differences between them therefore reflect
genuine estimand differences, not estimation noise.

\medskip\noindent\textbf{PLR-IV as a structural benchmark.}
PLR-IV is nonetheless informative.  It shares the residualized pseudo-IV structure with
Wald-AIPW \citep[cf.][]{knaus2024double} and admits an outcome-weight representation.
Including it allows one to assess the sensitivity of IV estimates to the constant-effect
restriction.  A large gap between PLR-IV and Wald-AIPW is direct evidence that effect
homogeneity is violated.

\begin{remark}
Both PLR-IV and Wald-AIPW use Double ML nuisance estimation---$\hat\ell,\hat r,\hat m$
estimated via cross-fitted machine learning---and both residualize $Y$, $D$, and $Z$ on
$X$.  Their difference lies solely in the moment condition: PLR-IV imposes
$\E[(\tY_i-\theta_0\tD_i)\tZ_i]=0$ under a constant-effect structural model, whereas
Wald-AIPW targets $\tau_{\mathrm{LATE}}$ through an augmented Wald score that allows
heterogeneous effects.
\end{remark}

\medskip\noindent\textbf{Reporting recommendation.}
PLR-IV results are reported as a separate \emph{structural DML-IV benchmark} rather than
as a LATE estimator.  Direct numerical comparisons to kappa or Wald-AIPW estimates
should acknowledge the distinct estimands.  If the two families yield similar estimates,
this is consistent with approximate effect homogeneity; if they diverge substantially,
the PLR-IV coefficient should not be interpreted as the complier average effect.

% ------------------------------------------------------------------
\bibliographystyle{apalike}
\bibliography{references}

\end{document}
